\documentclass[11pt]{article}
\usepackage{fontspec}
\usepackage{xeCJK}
\setCJKmainfont[ItalicFont=STSongti-SC-Bold]{Songti SC}
\usepackage{xurl}
\usepackage[a4paper,margin=25mm]{geometry}
\usepackage{amsmath,amssymb,amsthm,mathtools}
\usepackage{xcolor}
\usepackage{booktabs,array,enumitem,fancyhdr,hyperref}
\hypersetup{colorlinks=true,linkcolor=blue!50!black,urlcolor=blue!50!black}
\setlength{\emergencystretch}{2em}
\setlength{\parskip}{0.3em}
\linespread{1.16}
\allowdisplaybreaks
\newtheorem{theorem}{定理}[section]
\newtheorem{proposition}[theorem]{命题}
\newtheorem{lemma}[theorem]{引理}
\theoremstyle{definition}
\newtheorem{definition}[theorem]{定义}
\theoremstyle{remark}
\newtheorem{remark}[theorem]{说明}
\newcommand{\dd}{\,\mathrm d}
\newcommand{\ii}{\mathrm i}
\newcommand{\Rea}{\operatorname{Re}}
\newcommand{\Ima}{\operatorname{Im}}
\newcommand{\Li}{\operatorname{Li}}
\newcommand{\supp}{\operatorname{supp}}
\pagestyle{fancy}
\fancyhf{}
\fancyhead[L]{\small 学习笔记}
\fancyhead[R]{\small 初稿}
\fancyfoot[C]{\thepage}
\setlength{\headheight}{15pt}
\author{}
\date{2026年10月22日初稿}
\makeatletter
\renewcommand{\maketitle}{%
  \thispagestyle{plain}%
  \begin{center}
    {\LARGE\bfseries \@title\par}
    \vspace{0.6em}
    {\small \@date\par}
  \end{center}
  \vspace{0.6em}}
\makeatother
\newcommand{\draftnotice}{\noindent{\small
本文为供逐篇审阅的学习笔记初稿，不是新增研究成果。
文中区分直接证明与引用结果；正文不依赖本文。
本轮已作公式、条件及排版初步核查，尚待作者逐句审定。}\par\medskip}

\XeTeXlinebreaklocale "zh"
\renewcommand{\proofname}{证明}
\date{2026年10月24日初稿}
\title{黎曼猜想会不会不可判定}
\begin{document}
\maketitle

\noindent{\small 本篇讨论 RH 在形式系统中的逻辑地位。数理逻辑中的标准定理注明出处后引用；
其余命题给出证明。}\par\medskip

\section{问题}
自 Gödel 不完全性定理以来，一个自然的问题是：RH 是否可能在通常的公理系统
（如 Peano 算术 PA 或集合论 ZFC）中既不能证明也不能否证？
本篇说明，RH 的逻辑形式使这一问题具有特殊的答案：
\emph{若 RH 在 ZFC 中不可判定，则 RH 为真。}
这与连续统假设等真正的独立命题有本质区别。

\section{逻辑准备}
\subsection{算术分层}
以 $\Delta_0$ 记只含有界量词（$\forall x<t$，$\exists x<t$）的算术公式。称算术语句
\begin{itemize}
 \item 为 $\Sigma_1$ 型，若形如 $\exists n\,\varphi(n)$，$\varphi\in\Delta_0$；
 \item 为 $\Pi_1$ 型，若形如 $\forall n\,\varphi(n)$，$\varphi\in\Delta_0$。
\end{itemize}
$\Delta_0$ 公式在每个具体自然数上的真值可以机械判定。
由 Kleene 范式定理（或 MRDP 定理），每个可判定关系既可由 $\Sigma_1$ 公式、也可由 $\Pi_1$ 公式定义；
因此，若 $\varphi(n)$ 是任一可判定的性质，$\forall n\,\varphi(n)$ 也等价于一个 $\Pi_1$ 语句。
以下在此意义下使用 $\Pi_1$、$\Sigma_1$。

\subsection{$\Sigma_1$ 完全性}
\begin{theorem}[$\Sigma_1$ 完全性，引用]\label{thm:sigma1}
设 $T$ 是包含 Robinson 算术 Q 的理论（PA、ZFC 均是）。则每个在 $\mathbb N$ 中为真的 $\Sigma_1$ 语句都可在 $T$ 中证明。
\end{theorem}
其直观内容是：若 $\exists n\,\varphi(n)$ 为真，取一个具体的见证 $n_0$，
$\varphi(n_0)$ 的验证是有限计算，可逐步写成 $T$ 中的形式证明。
注意定理不要求 $T$ 是可靠的，只要求它包含 Q。

\section{RH 是 $\Pi_1$ 命题}
\begin{theorem}[Lagarias 2002，引用]\label{thm:lagarias}
记 $\sigma(n)=\sum_{d\mid n}d$，$H_n=\sum_{k\le n}1/k$。则 RH 成立当且仅当对一切 $n\ge1$，
\[
 \sigma(n)\le H_n+\exp(H_n)\log H_n .                    \tag{3.1}
\]
\end{theorem}
\begin{proposition}\label{prop:pi1}
RH 等价于一个 $\Pi_1$ 语句。
\end{proposition}
\begin{proof}
记 $R(n)=H_n+\exp(H_n)\log H_n$。对每个 $n$，$R(n)$ 是可计算实数：
存在算法，对输入 $(n,k)$ 输出有理数 $u_k(n)$，使 $R(n)\le u_k(n)\le R(n)+2^{-k}$
（$H_n$ 是有理数；$\exp$ 与 $\log$ 的 Taylor 展开有显式余项）。
由定理~\ref{thm:lagarias}，RH 的否定等价于存在 $n$ 使 $\sigma(n)>R(n)$。
由于不等号是严格的，这又等价于存在 $(n,k)$ 使 $\sigma(n)>u_k(n)$：
若 $\sigma(n)>R(n)$，取 $2^{-k}<\sigma(n)-R(n)$ 即可；反之由 $u_k(n)\ge R(n)$ 立得。
关系 $\sigma(n)>u_k(n)$ 是可判定的，所以 $\neg$RH 等价于一个 $\Sigma_1$ 语句，RH 等价于一个 $\Pi_1$ 语句。
\end{proof}

上述证明利用了 (3.1) 的否定是严格不等式这一点。
Davis、Matiyasevich 与 Robinson（1976）还给出了只含有理数运算的形式：
存在可计算函数 $g$，使 RH 等价于对一切正整数 $n$，
$\bigl(\sum_{k<g(n)}1/k-n^2/2\bigr)^2<36n^3$。
由 MRDP 定理，进一步可构造整系数多项式 $P$，使 RH 等价于 $P=0$ 没有自然数解。

也可直接从零点出发理解命题~\ref{prop:pi1}：若存在零点 $\rho_0$，$\beta_0\ne\tfrac12$，
则在一个包含 $\rho_0$、边界上无零点的有理矩形上，
辐角原理给出的零点个数与临界线上的零点个数之差可以由有限的区间运算严格算出。
因此“存在离线零点”可以由有限的见证证实。

\section{主要推论}
\begin{proposition}\label{prop:indep}
设 $T\supseteq$ Q。若 RH 在 $T$ 中不可否证，并且 $T$ 证明“RH $\Leftrightarrow\forall n\,\varphi(n)$”
（$\forall n\,\varphi(n)$ 为命题~\ref{prop:pi1} 中的 $\Pi_1$ 语句，且该等价在 $\mathbb N$ 中为真），则 RH 为真。
特别地，若 RH 独立于 ZFC，则 RH 为真。
\end{proposition}
\begin{proof}
若 RH 为假，则 $\exists n\,\neg\varphi(n)$ 是真的 $\Sigma_1$ 语句。
由定理~\ref{thm:sigma1}，$T\vdash\exists n\,\neg\varphi(n)$；
再由 $T$ 中的等价，$T\vdash\neg$RH，与假设矛盾。
ZFC 中可以形式化 Lagarias 定理的证明，故后一结论成立。
\end{proof}

由此得到两点。
\begin{enumerate}
 \item RH 的真值是确定的数学事实，不随集合论模型而改变。
 若存在离线零点，它是一个可以写下并验证的有限对象；若不存在，RH 即为真。
 不存在“在某些集合论宇宙中成立、在另一些中不成立”的情形。
 \item 若在某个元理论 $S$ 中证明了“RH 独立于 ZFC”，且 $S$ 足以形式化命题~\ref{prop:indep} 的论证，
 则 $S$ 也证明了 RH。因此，证明 RH 的独立性不会比证明 RH 本身更容易；
 “不可判定”不是绕开 RH 的出路。
\end{enumerate}

\section{与其他命题的比较}
\begin{itemize}
 \item \textbf{Goldbach 猜想}：“每个大于 $2$ 的偶数是两个素数之和”是 $\Pi_1$ 命题，
 与 RH 处于同一位置：若不可判定，则为真。
 \item \textbf{孪生素数猜想}：“对每个 $N$，存在 $p>N$ 使 $p,p+2$ 均为素数”是 $\Pi_2$ 命题。
 它的否定不再有有限见证，命题~\ref{prop:indep} 的论证不适用。
 \item \textbf{连续统假设}：不是算术命题。Gödel（1940）与 Cohen（1963）证明它独立于 ZFC（假定 ZFC 一致），
 并且其真值随集合论模型而改变。算术命题在 ZFC 的不同（有相同自然数的）模型之间真值不变，
 所以 RH 不可能以 CH 的方式独立。
 \item \textbf{已知的独立 $\Pi_1$ 命题}：由 Gödel 第二不完全性定理，
 若 ZFC 一致，则 $\mathrm{Con}(\mathrm{ZFC})$ 是真而不可证的 $\Pi_1$ 命题。
 因此“真而不可证的 $\Pi_1$ 命题”确实存在，RH 不可判定在逻辑上并非不可能。
 但目前已知独立于 ZFC 的 $\Pi_1$ 命题，或与一致性陈述直接相关，
 或由逻辑学家专门构造（例如 H.\ Friedman 的组合命题）；
 尚无经典数论猜想被证明属于此类。
\end{itemize}

\section{讨论}
\begin{enumerate}
 \item 命题~\ref{prop:indep} 是一个条件句：它不说明 RH 可判定，只说明“不可判定”蕴涵“为真”。
 是否存在 RH 在 ZFC 中不可证的可能，目前没有任何正面或反面的证据。
 \item 这一结论与学习笔记第 12 篇的不对称性相对应：RH 的否定若成立，可以由有限计算证实；
 RH 本身则须由证明确立。逐段验证 $\mathrm{RH}(T)$ 相当于在有界范围内搜索 $\Sigma_1$ 见证。
 \item 命题~\ref{prop:pi1} 依赖于 RH 的某个可判定实例形式（如 Lagarias 判据）。
 并非每个 RH 等价命题在字面上都是 $\Pi_1$ 的（例如 $M(x)=O(x^{1/2+\varepsilon})$ 对一切 $\varepsilon>0$），
 但它们经由定理彼此等价，因而逻辑地位相同。学习笔记第 14 篇讨论这些等价形式。
\end{enumerate}

\section*{参考}
\begin{itemize}
 \item J.\ C.\ Lagarias, An elementary problem equivalent to the Riemann hypothesis, \emph{Amer.\ Math.\ Monthly} 109 (2002), 534--543.
 \item M.\ Davis, Yu.\ Matijasevič, J.\ Robinson, Hilbert's tenth problem: Diophantine equations: positive aspects of a negative solution,
 \emph{Proc.\ Sympos.\ Pure Math.} 28 (1976), 323--378.
 \item P.\ Hájek, P.\ Pudlák, \emph{Metamathematics of First-Order Arithmetic}, Springer, 1993（$\Sigma_1$ 完全性）。
 \item K.\ Gödel, \emph{The Consistency of the Continuum Hypothesis}, 1940；P.\ J.\ Cohen, The independence of the continuum hypothesis, \emph{PNAS} 50 (1963).
\end{itemize}
\end{document}
